\documentclass[a4paper,11pt]{article}

% Package imports
\usepackage{amsmath, amssymb, amsthm}
\usepackage{graphicx}
\usepackage{fancyhdr}
\usepackage{geometry}
\usepackage{titlesec}
\usepackage{bm}
\usepackage[nolist]{acronym}
\usepackage[most]{tcolorbox}

\usepackage{needspace} % ensures rule+heading stay together

\newcommand{\rsectionpre}{%
  % Only if not at the top of a page
  \ifdim\pagetotal>0pt
    % If there isn't room for the rule + heading, break first
    \Needspace*{3\baselineskip}%
    % If we didn't break, draw the rule
    \ifdim\pagetotal>0pt
      \noindent\rule{\linewidth}{0.4pt}\par
      \vspace{2em}%
      \nobreak % avoid a break between rule and heading
    \fi
  \fi
}

\newcommand{\rsection}{\section}

% Tell titlesec how to format \ruledsection
\titleformat{\rsection}[block]
  {\rsectionpre\normalfont\Large\bfseries} % format (rule first)
  {\thesection}{1em}{}

% Optional: starred version \ruledsection* (unnumbered)
\titleformat*{\rsection}{\rsectionpre\normalfont\Large\bfseries}

% Page layout
\geometry{left=2cm, right=2cm, top=2.5cm, bottom=2.5cm}
\graphicspath{ {./images/} }

% Header definition
\pagestyle{fancy}
\fancyhf{}
\lhead{Reinforcement Learning - SS25}
\rhead{\thepage}

% Theorem and definition styles with rounded borders and numbering
\newtcolorbox{definitionbox}{colback=white, colframe=black, arc=5pt, boxrule=1pt, left=5pt, right=5pt, top=5pt, bottom=5pt}
\newtcolorbox{theorembox}{colback=white, colframe=black, arc=5pt, boxrule=1pt, left=5pt, right=5pt, top=5pt, bottom=5pt}

% Custom commands for numbered definitions and theorems
\newcounter{definitioncounter}
\newcounter{theoremcounter}

\setlength{\parindent}{0pt}
\setlength{\skip\footins}{0.5cm} % Adjust this to change the space above the footnotes
\raggedbottom

\newcommand{\definition}[1]{%
    \refstepcounter{definitioncounter}%
    \begin{definitionbox}
    \textbf{Definition \thedefinitioncounter:}
    \vspace{0.2cm}\newline
    #1
    \end{definitionbox}
    \vspace{0.3cm}
}

\newcommand{\theorem}[1]{%
    \refstepcounter{theoremcounter}%
    \begin{theorembox}
    \textbf{Theorem \thetheoremcounter:}
    \vspace{0.2cm}\newline
    #1
    \end{theorembox}
    \vspace{0.3cm}
}

\newcommand{\cf}[1]{\[#1\]}
\newcommand{\f}[1]{${#1}$}
\renewcommand{\b}[1]{\textbf{#1}}
\renewcommand{\it}[1]{\textit{#1}}
\newcommand{\pL}{\f{\mathcal{L}}}
\newcommand{\fL}{\mathcal{L}}

\begin{document}

% Title section
\begin{center}
    {\huge Reinforcement Learning - SS25 \par}
    \vspace{0.5cm}
    {\large Niklas Rodenbüsch \par}
    \vspace{0.5cm}
    {\large September 7, 2025 \par}
\end{center}
\vspace{0.5cm}



\section{Reinforcement Learning Fundamentals}

\definition{
    \textbf{Core RL Components:}
    \begin{itemize}
        \item \textbf{States} \f{s \in S}: Finite set of environment states
        \item \textbf{Actions} \f{a \in A}: Finite set of possible actions
        \item \textbf{Rewards} \f{r \in R}: Scalar feedback signals
        \item \textbf{Policies} \f{\pi}: Action selection strategies
        \begin{itemize}
            \item Deterministic: \f{\pi(s) = a}
            \item Stochastic: \f{\pi(a|s) = p(a|s)}
        \end{itemize}
        \item \textbf{Value Function} \f{v^\pi(s) = \mathbb{E}[G_t|s_t = s]}: Expected return from state
        \item \textbf{Action-Value Function} \f{q^\pi(s,a) = \mathbb{E}[G_t|s_t = s, a_t = a]}: Expected return from state-action pair
    \end{itemize}
}

The goal in RL is to maximize expected cumulative reward. We often use \textbf{discounted rewards} to handle uncertainty about future rewards:
\cf{G_t = r_{t+1} + \gamma r_{t+2} + \gamma^2 r_{t+3} + \cdots = \sum_{k=0}^{\infty} \gamma^k r_{t+k+1}}
where \f{\gamma \in [0,1]} is the discount factor.

\b{Bellman Equations} provide recursive relationships for value functions:
\cf{v^\pi(s) = \sum_a \pi(a|s) \sum_{s',r} p(s',r|s,a)[r + \gamma v^\pi(s')]}
\cf{q^\pi(s,a) = \sum_{s',r} p(s',r|s,a)[r + \gamma \sum_{a'} \pi(a'|s')q^\pi(s',a')]}

\section{Markov Decision Process (MDP)}

\definition{
    An MDP satisfies the \textbf{Markov Property}: the future depends only on the current state, not the history.
    \cf{p(s',r|s_t,a_t) = p(s',r|s_t,a_t,\ldots,s_0,a_0)}
}

\b{Example - Robot Cleaning}: A robot has charge levels \{high, low, none\} and actions \{explore, recharge\}. The transition probabilities depend only on current charge and action, demonstrating the Markov property.

\textbf{Optimal Policy} \f{\pi^*} satisfies: \f{v^{\pi^*}(s) \geq v^\pi(s)} for all states and policies.

\section{Multi-Armed Bandits (MAB)}

MAB problems have multiple actions but no states. The goal is to minimize \textbf{regret}:
\cf{\text{Regret} = \mathbb{E}\left[\sum_{t=1}^T (q^*(a^*) - q^*(a_t))\right]}

\subsection{Solution Methods:}

\textbf{$\epsilon$-Greedy}: Choose best action with probability \f{1-\varepsilon}, random with \f{\varepsilon}.
Action-value update: \f{q_{n+1} = q_n + \alpha(r_n - q_n)}

\textbf{UCB (Upper Confidence Bound)}: Select action based on optimism in face of uncertainty:
\cf{a_t = \arg\max_a \left[q_t(a) + c\sqrt{\frac{\ln t}{N_t(a)}}\right]}
The second term promotes exploration of less-tried actions.

\textbf{Softmax Selection}: Choose actions probabilistically based on their values:
\cf{p(a_t = a) = \frac{\exp(q_t(a)/\tau)}{\sum_{a'} \exp(q_t(a')/\tau)}}
where \f{\tau} controls exploration (\f{\tau \to 0}: greedy, \f{\tau \to \infty}: random).

\subsection{MAB Action Value Updates}
Incremental update rule for action values:
$$q_{n+1}(a) = q_n(a) + \frac{1}{n}[r_n - q_n(a)]$$

Or with step size $\alpha$:
$$q_{n+1}(a) = q_n(a) + \alpha[r_n - q_n(a)]$$

Where $n$ is the number of times action $a$ has been selected.

\section{Dynamic Programming}

Dynamic programming finds optimal policies using the Bellman equations when the MDP model is known.

\b{Policy Iteration}:
\begin{enumerate}
    \item \textbf{Policy Evaluation}: Compute \f{v^\pi} for current policy
    \item \textbf{Policy Improvement}: Update policy greedily: \f{\pi'(s) = \arg\max_a \sum_{s',r} p(s',r|s,a)[r + \gamma v^\pi(s')]}
    \item Repeat until convergence
\end{enumerate}

\b{Value Iteration}: Directly compute optimal value function:
\cf{v_{k+1}(s) = \max_a \sum_{s',r} p(s',r|s,a)[r + \gamma v_k(s')]}
Then extract optimal policy: \f{\pi(s) = \arg\max_a \sum_{s',r} p(s',r|s,a)[r + \gamma v^*(s')]}

\subsection{DP Method Characteristics}
\textbf{Policy Iteration:}\\
- Uses Bellman Expectation Equation for policy evaluation\\
- Guaranteed convergence to optimal policy\\
- Requires full model knowledge\\

\textbf{Value Iteration:}\\
- Uses Bellman Optimality Equation\\
- No explicit policy during iterations\\
- Extracts policy only at the end: $\pi(s) = \arg\max_a \sum_{s',r} p(s',r|s,a)[r + \gamma v^*(s')]$

\section{Monte Carlo Methods}

MC methods are \textbf{model-free} - they learn from experience without knowing the MDP model.

\b{Key Characteristics}:
\begin{itemize}
    \item Update only after complete episodes
    \item Low bias (use actual returns)
    \item High variance (single sample estimates)
    \item No bootstrapping required
\end{itemize}

\textbf{MC Prediction}: Estimate value function by averaging returns:
\cf{v^\pi(s) = \frac{1}{N(s)} \sum_{i: s_t^{(i)} = s} G_t^{(i)}}

\b{MC Control Methods}:
\begin{itemize}
    \item \textbf{On-policy}: Use same policy for behavior and learning (e.g., e-greedy)
    \item \textbf{Off-policy}: Learn target policy from behavior policy using importance sampling
\end{itemize}

\textbf{Importance Sampling} for off-policy learning:
\cf{\rho_{t:T-1} = \prod_{k=t}^{T-1} \frac{\pi(a_k|s_k)}{b(a_k|s_k)}}
This ratio corrects for the difference between target and behavior policies.

\section{Temporal Difference Learning}

TD methods combine advantages of DP (bootstrapping) and MC (model-free learning).

\b{Key Advantages}:
\begin{itemize}
    \item Online learning (update after each step)
    \item Lower variance than MC
    \item Can handle infinite episodes
\end{itemize}

\textbf{TD(0) Update}:
\cf{v(s_t) \leftarrow v(s_t) + \alpha[r_{t+1} + \gamma v(s_{t+1}) - v(s_t)]}

The \textbf{TD error} \f{\delta_t = r_{t+1} + \gamma v(s_{t+1}) - v(s_t)} drives learning.

\b{TD Control Methods}:

\textbf{SARSA (On-policy)}:
\cf{q(s_t,a_t) \leftarrow q(s_t,a_t) + \alpha[r_{t+1} + \gamma q(s_{t+1},a_{t+1}) - q(s_t,a_t)]}

\textbf{Q-Learning (Off-policy)}:
\cf{q(s_t,a_t) \leftarrow q(s_t,a_t) + \alpha[r_{t+1} + \gamma \max_a q(s_{t+1},a) - q(s_t,a_t)]}

\textbf{Expected SARSA}:
\cf{q(s_t,a_t) \leftarrow q(s_t,a_t) + \alpha[r_{t+1} + \gamma \sum_a \pi(a|s_{t+1})q(s_{t+1},a) - q(s_t,a_t)]}

\textbf{Double Q-Learning} addresses overestimation bias by maintaining two Q-functions:
\cf{q_1(s_t,a_t) \leftarrow q_1(s_t,a_t) + \alpha[r_{t+1} + \gamma q_2(s_{t+1}, \arg\max_a q_1(s_{t+1},a)) - q_1(s_t,a_t)]}

\subsection{Monte Carlo vs Temporal Difference Comparison}
\begin{tabular}{|l|l|l|}
\hline
Property & Monte Carlo & Temporal Difference \\
\hline
Variance & Higher & Lower \\
Bias & Unbiased (uses actual returns) & Biased (bootstrapping) \\
Initial value sensitivity & Less sensitive & More sensitive \\
Learning frequency & End of episode & Every time step \\
Episode requirement & Complete episodes & Can handle incomplete \\
\hline
\end{tabular}

\section{Function Approximation}

When state spaces are large or continuous, we need function approximation instead of tabular methods.

\definition{
    \textbf{Function Approximation}: Represent value functions as parameterized functions \f{\hat{v}(s,\mathbf{w})} where \f{\mathbf{w}} are learnable parameters.
}

\b{Linear Methods}: \f{\hat{v}(s,\mathbf{w}) = \mathbf{w}^T \mathbf{x}(s)}
where \f{\mathbf{x}(s)} is a feature vector. Linear TD methods have convergence guarantees.

\textbf{Semi-gradient TD}: 
\cf{\mathbf{w}_{t+1} = \mathbf{w}_t + \alpha[r_{t+1} + \gamma \hat{v}(s_{t+1},\mathbf{w}_t) - \hat{v}(s_t,\mathbf{w}_t)]\nabla \hat{v}(s_t,\mathbf{w}_t)}

Called "semi-gradient" because the target \f{r_{t+1} + \gamma \hat{v}(s_{t+1},\mathbf{w}_t)} also depends on \f{\mathbf{w}_t}.

\b{Deep Q-Networks (DQN)}:
\begin{itemize}
    \item Use neural networks to approximate Q-function
    \item \textbf{Experience Replay}: Store transitions, sample mini-batches randomly
    \item \textbf{Target Network}: Separate network for computing targets, updated periodically
\end{itemize}

\textbf{Deadly Triad}: Combining function approximation + bootstrapping + off-policy learning can cause instability.\\
\textbf{Solutions:}\\
- Target Networks (DQN): Separate network for computing targets\\
- Experience Replay: Break correlation in training data\\
- Double Q-Learning: Address overestimation bias

\section{Policy Gradient Methods}

Instead of learning value functions, directly optimize parameterized policies \f{\pi(a|s,\boldsymbol{\theta})}.

\theorem{
    \textbf{Policy Gradient Theorem}: The gradient of expected return is:
    \cf{\nabla J(\boldsymbol{\theta}) = \mathbb{E}_\pi[q^\pi(s,a) \nabla \ln \pi(a|s,\boldsymbol{\theta})]}
}

\textbf{REINFORCE}: Monte Carlo policy gradient method:
\cf{\boldsymbol{\theta}_{t+1} = \boldsymbol{\theta}_t + \alpha G_t \nabla \ln \pi(a_t|s_t,\boldsymbol{\theta}_t)}

High variance can be reduced using a \textbf{baseline} \f{b(s)}:
\cf{\boldsymbol{\theta}_{t+1} = \boldsymbol{\theta}_t + \alpha (G_t - b(s_t)) \nabla \ln \pi(a_t|s_t,\boldsymbol{\theta}_t)}

\b{Actor-Critic Methods}: Combine policy gradient (actor) with value function approximation (critic):
\begin{itemize}
    \item \textbf{Actor}: Policy \f{\pi(a|s,\boldsymbol{\theta})}
    \item \textbf{Critic}: Value function \f{\hat{v}(s,\mathbf{w})}
\end{itemize}
Update using TD error: \f{\delta_t = r_{t+1} + \gamma \hat{v}(s_{t+1},\mathbf{w}_t) - \hat{v}(s_t,\mathbf{w}_t)}

\subsection{Score Function}
For a policy $\pi(a|s,\theta)$, the score function is:
$$\nabla_\theta \ln \pi(a|s,\theta)$$

This appears in the policy gradient theorem and REINFORCE updates.

\subsection{REINFORCE with Baseline}
To reduce variance, subtract a baseline $b(s)$:
$$\theta_{t+1} = \theta_t + \alpha(G_t - b(s_t))\nabla \ln \pi(a_t|s_t, \theta_t)$$

Common baselines:
- State value function: $b(s) = v^{\pi}(s)$
- Average return
- Any function that doesn't depend on the action

% Add new section for DDPG
\subsection{Deep Deterministic Policy Gradient (DDPG)}
\textbf{Classification:} Off-policy method

\textbf{Key features:}
- Uses experience replay buffer
- Deterministic policy
- Actor-critic architecture
- Target networks for both actor and critic
- Gaussian exploration noise added to actions

\textbf{Why off-policy:}
- Experience replay uses transitions from previous policies
- State distribution not generated by current target policy

% Add Model-Based RL section
\section{Model-Based Reinforcement Learning}

\subsection{Overview}
Learn a model $f$ of the environment transition function $p(s',r|s,a)$, then use this model for planning.

\subsection{Dyna Algorithm}
Combines:
- Direct RL: Learn from real experience
- Planning: Learn from simulated experience using the model

\textbf{Dyna-Q Steps:}
\begin{enumerate}
\item Take action, observe reward and next state
\item Update Q-function with real experience
\item Update model with real experience  
\item Repeat n times: sample from model and update Q-function
\end{enumerate}

\subsection{Model Learning}
For Gaussian model: $f(s'|s,a;\theta) = \mathcal{N}(\mu(s,a;\theta), \Sigma)$

Where $\mu(s,a;\theta) = \theta^T \phi(s,a)$ with feature map $\phi(s,a)$.

\subsection{Advantages and Disadvantages}
\textbf{Advantages:}
- More sample efficient
- Can plan ahead
- Can simulate experience

\textbf{Disadvantages:}
- Model errors compound
- Computationally expensive
- Model bias affects performance

\subsection{Planning Methods}
\textbf{Full Dynamic Programming:}
- Use learned model in Bellman equations
- Computationally expensive for large state spaces

\textbf{Sampling-based Planning:}
- Sample transitions from learned model
- Apply MC or TD methods to sampled experience
- More computationally efficient but higher variance

\section{Advanced Topics}

\b{n-step Methods}: Bridge between TD ($n=1$) and MC ($n=\infty$):
\cf{G_{t:t+n} = r_{t+1} + \gamma r_{t+2} + \cdots + \gamma^{n-1} r_{t+n} + \gamma^n \hat{v}(s_{t+n})}

\textbf{TD($\lambda$)}: Combines multiple n-step returns using eligibility traces:
\cf{G_t^\lambda = (1-\lambda) \sum_{n=1}^{\infty} \lambda^{n-1} G_{t:t+n}}

\textbf{Eligibility Traces}: Track which states should receive credit for rewards:
\cf{\mathbf{z}_t = \gamma \lambda \mathbf{z}_{t-1} + \nabla \hat{v}(s_t,\mathbf{w}_t)}

\b{Model-Based RL}:
\begin{itemize}
    \item Learn environment model \f{p(s',r|s,a)}
    \item Use model for planning (generating simulated experience)
    \item \textbf{Dyna}: Interleave real experience with planning updates
\end{itemize}

\textbf{Monte Carlo Tree Search (MCTS)}:
\begin{enumerate}
    \item \textbf{Selection}: Navigate tree using tree policy
    \item \textbf{Expansion}: Add new nodes
    \item \textbf{Simulation}: Random rollout from leaf
    \item \textbf{Backpropagation}: Update values along path
\end{enumerate}

\b{Key Tradeoffs}:
\begin{itemize}
    \item \textbf{Exploration vs Exploitation}: Balance trying new actions vs using known good ones
    \item \textbf{Bias vs Variance}: TD (higher bias, lower variance) vs MC (lower bias, higher variance)  
    \item \textbf{Sample Efficiency vs Stability}: Off-policy (more efficient) vs on-policy (more stable)
\end{itemize}

\section{Important Formulations for Exam}

\textbf{UCB Confidence Bound:}
$$\sqrt{\frac{\ln t}{N_t(a)}}$$

\textbf{TD Error:}
$$\delta_t = r_{t+1} + \gamma v(s_{t+1}) - v(s_t)$$

\textbf{Importance Sampling Ratio:}
$$\rho_{t:T-1} = \prod_{k=t}^{T-1} \frac{\pi(a_k|s_k)}{b(a_k|s_k)}$$

\textbf{Semi-gradient Update:}
$$w_{t+1} = w_t + \alpha[r_{t+1} + \gamma \hat{v}(s_{t+1}, w_t) - \hat{v}(s_t, w_t)]\nabla \hat{v}(s_t, w_t)$$

% Add problem formalization tips
\subsection{MDP Formalization Tips}
When formalizing a problem as an MDP:
\begin{enumerate}
\item Clearly define state space $S$
\item Define action space $A$ (may be state-dependent)
\item Specify reward function $R$
\item Define transition probabilities $p(s',r|s,a)$
\item Only include transitions with non-zero probability
\end{enumerate}

% Add regret analysis
\subsection{Regret Analysis}
\textbf{Definition:}
$$\text{Regret} = E\left[\sum_{t=1}^T (q^*(a^*) - q^*(a_t))\right]$$

\textbf{Asymptotic behavior:}
- $\epsilon$-greedy: Linear regret (always explores with fixed probability)
- UCB: Logarithmic regret (exploration decreases over time)


\end{document}